\section{Foundations, Model Families, and System Building Blocks}
\label{sec:foundations}

The capability chapters focus on medical implementations. This section instead isolates the technical and conceptual building blocks that those systems borrow from general world models and adjacent decision-learning methods. The separation prevents two common distortions. General agents and driving simulators should not inflate estimates of medical capability, and a clinical policy should not be called model-based merely because it acts over time. These works are used as comparators: they clarify what state, dynamics, interaction, and planning can mean before those functions are subjected to clinical evidence requirements.

\subsection{General World Models and Latent Dynamics}

At minimum, a world model separates a represented state from a rule for changing that state. Transformer world models show that both the representation and transition can be learned as sequence prediction, with performance depending on how observations are tokenized, compressed, and rolled forward \citep{Micheli2023TransformersAreSample,Robine2023TransformerBasedWorld}. This abstraction transfers naturally to EHR events, longitudinal images, and physiological trajectories. Its medical interpretation does not transfer automatically. Tokens in a game describe an environment whose rules are defined; tokens in an EHR combine disease, observation, and clinician behavior.

The main architectural lesson is consequently functional. A latent state should retain information needed for the claimed transition and outcome; the transition should expose its temporal target and any action pathway; and the decoder or observation model should not be mistaken for the underlying dynamics. A larger latent model may improve generative fit while leaving state observability, patient identity, and action support unresolved. The later taxonomy therefore classifies the implemented contract rather than the capacity of the backbone.

\subsection{Predictive Representations and Generative Environments}

Predictive representation learning offers a route to world-model state without requiring pixel-perfect reconstruction. Joint-embedding objectives and video feature prediction learn representations by predicting missing or future latent content, while generative-environment work turns learned visual regularities into interactive sequences \citep{Assran2023Selfsupervisedlearning,Bardes2024RevisitingFeaturePrediction,Bruce2024Genie}. Hybrid use of predictive representations and diffusion priors further illustrates that representation and generation can be separate modules \citep{Park2026BeyondGenerativePriors}. These distinctions matter in medicine because a model can learn anatomically useful state while remaining below the existence boundary for a dynamic world model.

Generative realism introduces a parallel ambiguity. A model can produce a coherent frame by learning appearance statistics, but a clinical transition requires persistence of patient-specific anatomy, pathology, and physiology. Interactive generation adds another test: changing an input should alter the future through a stable action interface rather than through unconstrained prompting. These requirements motivate the Level~1/Level~2 distinction between state and directed future, and the Level~2/Level~3 distinction between temporal conditioning and executable action.

\subsection{Planning, Control, and Imagined Rollouts}

The clearest general-world-model lineage couples learned dynamics to action selection. Early latent world models and imagination-based control systems established the pattern of encoding observations, rolling latent state forward, and training or evaluating a controller inside the learned environment \citep{unknown2018WorldModels,Hafner2019DreamControlLearning,Hafner2019LearningLatentDynamics}. Dreamer variants extend this pattern across discrete and diverse domains, while MuZero demonstrates planning with a learned model that need not reconstruct every environmental detail \citep{Hafner2021MasteringAtariDiscrete,Hafner2023DreamerV3,Schrittwieser2020MuZero}. The transferable principle is that a model is decision-relevant when its predictions preserve what the planner needs, not when every observation is reproduced faithfully.

Model-predictive and hierarchical approaches place stronger emphasis on short-horizon replanning, task structure, and physical interaction. DayDreamer brings learned world models to robot learning; TD-MPC and TD-MPC2 optimize action sequences through learned latent dynamics; Director introduces hierarchical planning from pixels \citep{Wu2022DaydreamerWorldModels,Hansen2022TemporalDifferenceLearning,Hansen2024TdMpc2Scalable,Hafner2022DeepHierarchicalPlanning}. These designs offer useful engineering patterns for acquisition guidance, surgical control, and physiological management. They also reveal a medical risk: optimization searches for trajectories that exploit whatever the learned model and objective fail to constrain.

Recent systems scale the environment model itself. Diffusion-based world modeling, predictive video models, physical-AI platforms, autonomous-driving world models, and interactive real-world simulators expand visual fidelity, controllability, and domain coverage \citep{Alonso2024DiffusionWorldModeling,Assran2025VJepa2,NVIDIA2025Cosmos,Hu2023Gaia1Generative,Russell2025Gaia2Controllable,Yang2024LearningInteractiveReal}. They are useful comparators for surgical video, medical acquisition, and embodied simulation, but they remain outside the medical implementation denominator. Their actions, rewards, and safety constraints are defined for other environments. Medical transfer requires re-establishing the state, action, outcome, and validation contract rather than assuming that scale supplies clinical semantics.

\subsection{Clinical Decision Learning without an Explicit World Model}

Clinical reinforcement and imitation learning provide an important negative boundary. Policies have been learned for ventilator weaning and sepsis treatment, including approaches designed to reduce distribution-shift risk, while imitation learning has been used for surgical task execution \citep{Prasad2017ReinforcementLearningApproach,Raghu2017ContinuousStateSpace,Kaushik2022ConservativeQLearning,Kim2024SurgicalRobotTransformer}. These systems can be sequential and action-oriented without containing a world model. A value function, direct policy, or behavior-cloning transformer does not become model-based unless an explicit learned or mechanistic transition is used to evaluate, train, or optimize actions.

This boundary prevents application labels from deciding capability. ``Treatment optimization,'' ``control,'' and ``robotic autonomy'' describe intended use, not internal model use. Conversely, a medical simulator can satisfy the model component while supporting only training or offline evaluation rather than bedside control. The Level~5 contract therefore asks where the model enters the decision loop and whether removing or replacing it changes action ranking, policy learning, or control.

\subsection{Mechanistic, Causal, and Hybrid Building Blocks}

Medical world models draw on two additional traditions that are less prominent in general benchmarks. Mechanistic models specify physiological, anatomical, pharmacological, or physical structure and can expose intervention pathways directly. Causal models specify the estimand and assumptions required to compare outcomes under alternative actions. Sequence-based counterfactual estimators, G-computation networks, and causal transformers show how time-varying treatment and outcome models can be made explicit \citep{Bica2020CRN,Li2020GNet,Melnychuk2022CausalTransformer,Seedat2022TECDE}. These methods are discussed in depth at Level~4 because they implement counterfactual task contracts, even when they are not branded as world models.

Hybrid systems can combine learned representations with mechanistic state, structural constraints, or causal outcome heads. The potential advantage is not that hybrid models are valid by construction. It is that each component can carry a named responsibility: perception, state estimation, transition, intervention response, uncertainty, or planning. This modularity makes assumptions and failure modes more inspectable. It also exposes interface errors, such as a patient-specific encoder feeding population-level dynamics or a planner treating an uncertain mechanistic parameter as known.

\subsection{Cellular World Models and Biomedical Boundaries}

World-model language is also moving below the patient scale. AlphaCell, Lingshu-Cell, and GeneJEPA model cellular or transcriptomic state and, in some cases, perturbation-induced dynamics \citep{Chuai2026BuildingWorldModel,Zhang2026LingshuCellGenerative,Litman2025GenejepaPredictiveWorld}. These systems are scientifically relevant because they sharpen the same questions about state, transition, perturbation, and prediction. They are not counted as medical patient-level implementations in the primary capability distribution, since their states, actions, and validation targets operate at a different biological scale.

The boundary is analytical rather than dismissive. Cellular models may eventually supply mechanistic priors or intervention-response components for multiscale patient models. Population-health simulators may likewise inform policy without representing an individual patient. Keeping these scales visible allows the survey to discuss transfer and integration without conflating cellular perturbation, patient treatment, and health-system policy as equivalent actions.

\subsection{What Transfers to Medicine}

Across these foundations, four lessons recur. First, predictive representation is a prerequisite but not a transition. Second, a generative transition is not necessarily an action-response model. Third, action selection does not imply model-based planning. Fourth, planning through a model does not establish causal or clinical validity. These separations define the lower and upper boundaries of the five capability levels.

The medical opportunity lies in combining components that currently live in different traditions: multimodal state estimation, temporal and mechanistic dynamics, explicit interventions, counterfactual identification, uncertainty-aware planning, and decision-aligned validation. The analytical framework in the next section turns this component view into a claim-to-evidence contract.
